\chapter{Guia de Sintaxe e Modelagem de Macros Lógicas}
\label{ap:guia-sintaxe-macros}

Este apêndice serve como guia de referência técnico para a especificação do formato Twee 3 e a sintaxe de macros lógicas SugarCube suportadas pelo editor do \gls{piap}. Inclui exemplos práticos de modelação de mecânicas de jogo comuns em narrativas interativas.

\section{Estrutura de um Documento Twee 3}
O formato Twee 3 é uma representação textual de histórias não-lineares estruturada sob a forma de passagens de texto delimitadas por cabeçalhos iniciados por duas colunas (\texttt{::}). A estrutura de um projeto compila-se em passagens de sistema obrigatórias e passagens narrativas (cenas).

\subsection{Passagens Especiais de Sistema}
Para ser considerado válido, um ficheiro Twee 3 importado no \gls{piap} deve conter os seguintes nós de sistema:
\begin{itemize}
    \item \textbf{StoryTitle:} Contém exclusivamente uma única linha de texto plano com o título geral do projeto.
    \item \textbf{StoryData:} Contém um bloco estruturado em formato \gls{json} com as configurações globais da história. Os campos obrigatórios são o \texttt{ifid} (um identificador único de 128 bits para o projeto), o formato (\texttt{SugarCube}), e o nó de arranque (\texttt{start}):
    \begin{code}
    :: StoryData
    {
      "ifid": "A4B7D9E0-C2B1-4F2A-8E6D-009988776655",
      "format": "SugarCube",
      "start": "cena_inicial",
      "zoom": 1.0
    }
    \end{code}
    
    \item \textbf{StoryInit:} O nó lógico de inicialização. É executado uma única vez ao carregar a história na memória, antes de renderizar qualquer cena ao utilizador. Destina-se à declaração de variáveis globais e valores por defeito:
    \begin{code}
    :: StoryInit
    <<set $tem_chave = false>>
    <<set $moedas = 0>>
    <<set $nome_jogador = "Tomas">>
    \end{code}
\end{itemize}

\subsection{Passagens Narrativas (Cenas)}
As cenas contêm o texto da história, as atribuições de variáveis em tempo de execução e as escolhas de navegação do jogador. O cabeçalho de uma passagem narrativa no \gls{piap} pode incluir metadados de posicionamento posicional do canvas de grafos (X e Y) e Zonas (parentId) expressos em \gls{json}:
\begin{code}
:: cena_inicial {"position":{"x":100,"y":150},"size":{"width":250,"height":120},"parentId":"capitulo1"} [tag1 tag2]
Aqui começa a aventura. O mercado está movimentado e podes ver uma taberna.
[[Entrar na taberna|cena_taberna]]
[[Ir para o beco escuro|cena_beco]]
\end{code}

\section{Sintaxe de Macros Lógicas (SugarCube)}
O compilador JIT e interpretador integrado no editor do \gls{piap} processa as seguintes macros lógicas para manipulação de variáveis e controlo de fluxo:

\subsection{Macro de Atribuição: \texttt{<<set>>}}
Utilizada para modificar o valor de uma variável. O \gls{piap} suporta atribuições diretas, aritméticas e de strings:
\begin{itemize}
    \item \textbf{Atribuição Booleana:} \texttt{<<set \$tem\_espada = true>>}
    \item \textbf{Incrementação/Decrementação:} \texttt{<<set \$vida = \$vida - 10>>} ou \texttt{<<set \$moedas = \$moedas + 5>>}
    \item \textbf{Atribuição de Texto:} \texttt{<<set \$armadura = "cota\_de\_malha">>}
\end{itemize}

\subsection{Macro de Controlo Condicional: \texttt{<<if>>}}
Permite criar ramificações dinâmicas na história com base nas variáveis lógicas. A estrutura suporta condicionamento encadeado:
\begin{code}
<<if $moedas >= 10>>
  Podes comprar a poção de cura.
  [[Comprar poção|cena_compra]]
<<else>>
  Não tens moedas suficientes para falar com o mercador.
  [[Voltar à praça|cena_praca]]
<<endif>>
\end{code}

\section{Padrões Comuns de Modelação Narrativa}
Apresentam-se de seguida soluções estruturadas para programar mecânicas clássicas de jogos de aventura e RPG no editor do \gls{piap}:

\subsection{Padrão 1: Inventário e Portas Trancadas}
Mecânica em que o acesso a um determinado nó está bloqueado até que o jogador recolha um item de chave noutro local.
\begin{enumerate}
    \item Inicializar a variável booleana no nó de metadados \texttt{StoryInit}:
    \begin{code}
    <<set $chave_portal = false>>
    \end{code}
    \item Na sala onde a chave é recolhida (\texttt{cena\_bau}), alterar o valor da variável:
    \begin{code}
    :: cena_bau
    Abres o baú de madeira e encontras uma chave dourada e brilhante.
    <<set $chave_portal = true>>
    [[Voltar ao salão|cena_salao]]
    \end{code}
    \item Na sala do portal (\texttt{cena\_salao}), condicionar o acesso à porta trancada:
    \begin{code}
    :: cena_salao
    Estás em frente ao portal de pedra. A fechadura parece requerer uma chave especial.
    [[Voltar ao bau|cena_bau]]
    <<if $chave_portal>>
      [[Utilizar a chave e passar o portal|cena_vitoria]]
    <<else>>
      O portal está firmemente trancado.
    <<endif>>
    \end{code}
\end{enumerate}

\subsection{Padrão 2: Contador de Tentativas e Fim de Jogo}
Mecânica em que o jogador tem um limite de tentativas ou pontos de vida antes de ser derrotado.
\begin{enumerate}
    \item Inicializar o contador numérico em \texttt{StoryInit}:
    \begin{code}
    <<set $tentativas = 3>>
    \end{code}
    \item Na cena onde ocorre o erro ou armadilha (\texttt{cena\_armadilha}), subtrair uma tentativa e verificar se o valor atinge zero:
    \begin{code}
    :: cena_armadilha
    O mecanismo da parede dispara flechas na tua direção! Conseguiste esquivar-te, mas perdeste energia.
    <<set $tentativas = $tentativas - 1>>
    <<if $tentativas > 0>>
      Tens ainda $tentativas pontos de vida restantes.
      [[Tentar novamente com cautela|cena_salao]]
    <<else>>
      As tuas forças esgotaram-se.
      [[Game Over|cena_derrota]]
    <<endif>>
    \end{code}
\end{enumerate}

\subsection{Padrão 3: Caminhos Ocultos e Nós Secretos}
Utilizado para cenas extra ou segredos da história recomendados apenas para leitores atentos.
\begin{itemize}
    \item Para evitar que nós em desenvolvimento ou caminhos secretos estraguem a compilação ou apareçam nos relatórios de erros de caminhos órfãos, adicione a etiqueta \texttt{secreto} nas tags do nó:
    \begin{code}
    :: cena_secreta {"position":{"x":500,"y":500}} [secreto]
    Encontraste a sala secreta do programador! Parabéns!
    [[Voltar à praça|cena_praca]]
    \end{code}
    O validador estático de grafos (\gls{bfs}) do \gls{piap} ignorará este nó ao verificar o alcance obrigatório da história, prevenindo falsos positivos de caminhos inacessíveis.
\end{itemize}